% !PRINTED 541
\weberpage{541}

%<<CONT
$F(x, \eta) = 0$ is, however, again cyclic in $\Omega(\eta)$, since the cyclic functions of its roots belong to this field.

The cyclic equations, like all \gesperrt{Abel}ian equations, have the property that every root is rationally expressible in terms of every other. Here these expressions can be arranged cyclically in the following manner. If $\alpha, \alpha_1, \alpha_2 \ldots \alpha_{m-1}$ are the roots of the cyclic equation $F(x) = 0$, then the polynomial of degree $(m-1)$ in $x$
\[
F(x)\left(\frac{\alpha_1}{x-\alpha} + \frac{\alpha_2}{x-\alpha_1} + \cdots + \frac{\alpha}{x-\alpha_{m-1}}\right) = \Psi(x)
\]
is unchanged by the permutation $\pi$ and therefore contained in $\Omega$. If we set $x = \alpha, \alpha_1 \ldots \alpha_{m-1}$ therein and introduce the sign
\[
\frac{\Psi(x)}{F'(x)} = \Theta(x)
\]
it follows that
\weq{11}{\alpha_1 = \Theta(\alpha), \alpha_2 = \Theta(\alpha_1), \ldots \alpha_{m-1} = \Theta(\alpha_{m-2}), \alpha = \Theta(\alpha_{m-1}),}
and this holds, whether $F(x)$ be reducible or irreducible, provided only that $F(x)$ and $F'(x)$ have no common divisor.

The solution of the cyclic equations of degree $m$ is hereby made dependent on the solution of cyclic equations whose degrees are the prime factors of $m$. If, for example, $m$ is a power of $2$, the solution is accomplished by a sequence of square roots. In this way \gesperrt{Gauss} first treated the cyclotomic equations\footnote{\gesperrt{Gauss}, Disquisitiones arithmeticae, Sectio VII.}.

We finally attach here the remark, important for what follows, that for a prime degree the concepts of the \gesperrt{normal equation and of the cyclic equation} coincide. For in any case a cyclic equation of prime degree, since it is irreducible, is a normal equation. And if, conversely, the degree $n$ of a normal equation, which is at the same time the degree of the group of this equation, is a prime, and $\pi$ is a non-identical permutation of this group, then the degree of $\pi$, which must indeed be a divisor of $n$, is equal to $n$, and the group of the equation is $1, \pi, \pi^2 \ldots \pi^{n-1}$, hence cyclic.
